\documentclass[twocolumn]{iopjournal}

\usepackage{amsmath,amssymb,amsfonts}

\fancyhfoffset[L]{0mm}
\setlength{\marginparwidth}{0pt}
\setlength{\marginparsep}{0pt}

% ==========================================
% ZENODO PREPRINT HEADERS (ACTIVE)
% ==========================================
\renewcommand{\articletype}[1]{%
  {\vspace*{-8mm}\noindent \Large \sf Preprint}\par
  \vspace*{6mm}\noindent{\scriptsize \sf{\bfseries \MakeUppercase{#1}}}}

\fancyhead[L]{}
\fancyhead[R]{}

\begin{document}

\articletype{Supplementary Material}

\title{Supplementary Material: The TetB Field Fraction\\
from Exchange Symmetry}

\author{Stephen Nelson}

\affil{Field Dynamics Research}

\ead{stephen@fielddynamics.org}

\begin{abstract}
The companion letter~\cite{Nelson2026a} proved that the tetrad
equation produces two bounded regions related by spatial parity:
TetA (the identity component) and TetB (the parity component).
This note derives the fraction of the total gravitational field
response carried by each sector. The exchange symmetry between
TetA and TetB, combined with rational bijectivity on the
projective line, uniquely determines the TetA fraction as
$\mu(x) = x/(x+1)$ and the TetB fraction as
$1 - \mu(x) = 1/(x+1)$, where $x = g/a_0$ is the
dimensionless acceleration. The derivation uses no field
equation, no coupling polynomial, and no equilibrium
postulates. It requires only the parity relation between the
two bounded regions of the tetrad equation together with the
rational-bijection classification on the projective line.
\end{abstract}

% ==================================================================

\section*{The two sectors}

The companion letter~\cite{Nelson2026a} proves that the
spatial sector of the tetrad equation
$g_{\mu\nu} = \eta_{ab}\,e^a{}_\mu\,e^b{}_\nu$ produces
exactly two bounded tetrahedral regions. Under spatial parity
$P: e^i \to -e^i$, the identity component gives TetA and the
parity component gives TetB. Their union is the stellated
octahedron. The tetrad equation is quadratic in the frame
field: TetA and TetB are its two roots in the parity
direction, and the metric
$g_{\mu\nu} = \eta_{ab}\,e^a{}_\mu\,e^b{}_\nu$ is the
quadratic form common to both.

Because TetA and TetB together constitute the two bounded
regions of the tetrad equation, the gravitational field
response at each point partitions between them. Let $\mu(x)$
denote the fraction of the total field response carried by
TetA at dimensionless acceleration $x = g/a_0$, so that
$1 - \mu(x)$ is the fraction carried by TetB. The two fractions
sum to
unity because $\mu$ measures the relative partition, not the
absolute field strength.

% ==================================================================

\section*{The exchange symmetry}

Spatial parity exchanges TetA and TetB, swapping the two
fractions: if $\mu$ is the TetA fraction before the exchange,
then $1 - \mu$ is the TetA fraction after it.

The natural variable for a two-sector partition is the odds
ratio between TetA and TetB:
\begin{equation}
  \Omega = \frac{\mu}{1 - \mu}
  = \frac{\text{TetA fraction}}{\text{TetB fraction}}.
  \label{eq:odds}
\end{equation}
Under exchange, the numerator and denominator swap, so
$\Omega \to 1/\Omega$ by construction. This inversion is
algebraic, requiring no assumption about how the acceleration
transforms.

Two endpoint conditions follow from the partition definitions:
\begin{itemize}
\item $\mu(0) = 0$: at zero acceleration, the field sits
  entirely in TetB. Equivalently, $\Omega(0) = 0$.
\item $\mu(x) \to 1$ as $x \to \infty$: at large
  acceleration, the field sits entirely in TetA. Equivalently,
  $\Omega \to \infty$.
\end{itemize}

The dimensionless acceleration $x = g/a_0$ is constructed from
the metric, which is determined by the 40~symmetric
(Christoffel) components of the connection; these are the TetA
sector~\cite{Nelson2026a}. The variable $x$ therefore measures
TetA-sector dominance: $x > 1$ where TetA dominates, $x < 1$
where TetB dominates, and $x = 1$ at the parity midpoint.
Since $x$ is intrinsically a TetA-sector quantity, exchange
inverts it to $1/x$; combined with $\Omega \to 1/\Omega$, the
full exchange symmetry requires $\Omega(1/x) = 1/\Omega(x)$.

% ==================================================================

\section*{Classification}

The map $\Omega(x)$ must be rational, to preserve the
algebraic structure of the tetrad equation, and bijective,
so that a single acceleration determines a unique partition.
Rationality is inherited from the polynomial character of the
tetrad equation itself; bijectivity is the physical requirement
that the field be uniquely decodable from the observed
acceleration.
A rational bijection on the projective line is necessarily
fractional linear~\cite{Hartshorne1977}:
\begin{equation}
  \Omega(x) = \frac{ax + b}{cx + d}.
  \label{eq:mobius}
\end{equation}

% ==================================================================

\section*{The unique solution}

\paragraph{Proposition.} \textit{Let
$\Omega(x) = (ax+b)/(cx+d)$ be a rational bijection
satisfying $\Omega(1/x) = 1/\Omega(x)$ and $\Omega(0) = 0$.
Then $\Omega(x) = x$.}

\paragraph{Proof.}
The endpoint condition $\Omega(0) = b/d = 0$ requires $b = 0$.
The map reduces to $\Omega(x) = ax/(cx + d)$.

Substituting into $\Omega(1/x) = 1/\Omega(x)$ and
cross-multiplying gives $a^2 x = (c + dx)(cx + d)$.
Comparing coefficients of $x^2$, $x^1$, and $x^0$:
\begin{align}
  x^2: \quad & cd = 0, \label{eq:x2} \\
  x^1: \quad & a^2 = c^2 + d^2, \label{eq:x1} \\
  x^0: \quad & cd = 0. \label{eq:x0}
\end{align}
If $d = 0$: $\Omega = a/c$, a constant, which is not
bijective. Rejected.

If $c = 0$: $a^2 = d^2$, so $a = \pm d$. The branch
$a = -d$ gives $\Omega(x) = -x$, which is incompatible
with $\Omega > 0$ for $x > 0$. Therefore $a = d$ and
\begin{equation}
  \Omega(x) = x.
  \label{eq:omega_result}
\end{equation}
\hfill$\square$

% ==================================================================

\section*{The TetA and TetB fractions}

From $\Omega(x) = x$, the TetA fraction follows by
normalisation:
\begin{equation}
  \mu(x) = \frac{\Omega}{1 + \Omega} = \frac{x}{x + 1}.
  \label{eq:mu}
\end{equation}
The TetB fraction is the complement:
\begin{equation}
  \boxed{1 - \mu(x) = \frac{1}{x + 1}.}
  \label{eq:tetb_result}
\end{equation}
This is the fraction of the total gravitational field
response carried by TetB, the parity component of the
tetrad equation, at dimensionless acceleration $x$.

At large acceleration ($x \gg 1$), the TetB fraction
vanishes: the field response is carried almost entirely by
TetA. At the parity midpoint ($x = 1$), TetA and TetB
contribute equally: $\mu = 1/2$, $1 - \mu = 1/2$. At low
acceleration ($x \ll 1$), the TetB fraction approaches
unity: the field response is carried almost entirely by TetB.

% ==================================================================

\section*{Consequences for the connection}

The companion letter~\cite{Nelson2026a} establishes that
TetA is the identity component and TetB is the parity
component under spatial parity $P: e^i \to -e^i$. The
affine connection $\Gamma^I{}_{JK}$ decomposes under index
exchange $(J \leftrightarrow K)$ into 40~symmetric components
(Christoffel) and 24~antisymmetric components (torsion).

A torsion-free connection is uniquely determined by the metric
and its first derivatives. The metric
$g_{\mu\nu} = \eta_{ab}\,e^a{}_\mu\,e^b{}_\nu$ is bilinear
in the frame field; under $P: e^i \to -e^i$, both factors
acquire a sign flip and the two signs cancel, leaving the
metric invariant under TetA$\,\leftrightarrow\,$TetB exchange.
A torsion-free connection, built entirely from this invariant
metric, therefore cannot distinguish TetA from TetB. This
root-blindness is a structural property of the quadratic map
$\phi: e \mapsto \eta(e,e)$: the image of a 2-to-1 map does
not resolve which pre-image produced it. TetA and TetB are
geometrically distinct regions~\cite{Nelson2026a}; their
indistinguishability under the torsion-free connection reflects
the character of $\phi$, not a physical equivalence of the two
roots. The
torsion sector is the necessary carrier of the TetA/TetB
distinction at the connection level.

General relativity's zero-torsion condition
$\Gamma^I{}_{[JK]} = 0$ removes this carrier from the
connection. In the partition language of this note,
this sets $1 - \mu = 0$ and therefore $\mu = 1$: the full
field response is attributed to a single effective sector.

The exchange symmetry proves that $\mu \neq 1$. The TetB
fraction is $1/(x+1)$, not zero. The field response that
general relativity removes from the connection is:
\begin{equation}
  \Delta\mu(x) = \frac{1}{x + 1}.
  \label{eq:delta_mu}
\end{equation}

% ==================================================================

\section*{The dark-to-baryon ratio}

If the TetB field response is interpreted as a mass
contribution (as the cold dark matter framework does), the
ratio of the removed response to the retained response is:
\begin{equation}
  \frac{\Delta\mu}{\mu}
  = \frac{1/(x+1)}{x/(x+1)}
  = \frac{1}{x}.
  \label{eq:ratio}
\end{equation}
This is an algebraic identity of the constitutive law. It
does not depend on any galaxy model, halo profile, or
observational data.

The cold dark matter framework decomposes the observed
gravitational field as $g_{\mathrm{obs}} = g_{\mathrm{baryon}}
+ g_{\mathrm{DM}}$. By definition,
$g_{\mathrm{baryon}} = \mu \cdot g_{\mathrm{obs}}$ and
$g_{\mathrm{DM}} = (1 - \mu) \cdot g_{\mathrm{obs}}$.
Therefore:
\begin{equation}
  \frac{g_{\mathrm{DM}}}{g_{\mathrm{baryon}}}
  = \frac{1 - \mu}{\mu}
  = \frac{1}{x}.
  \label{eq:cdm_ratio}
\end{equation}
The dark-to-baryon mass ratio that the cold dark matter
framework fits with a multi-parameter halo profile reduces
to $1/x$: the reciprocal of the dimensionless acceleration.
This follows from the exchange symmetry of the two
parity-related sectors of the tetrad equation, with no free
parameters and no halo model.

% ==================================================================

\section*{Summary}

The tetrad equation produces two bounded regions related by
spatial parity: TetA (identity) and TetB (parity). The
exchange symmetry, together with rational bijectivity, uniquely
determines the partition of the gravitational field response
between them:
\begin{align}
  \text{TetA:} \quad & \mu(x) = \frac{x}{x + 1},
  \label{eq:sum_mu} \\[4pt]
  \text{TetB:} \quad & 1 - \mu(x) = \frac{1}{x + 1}.
  \label{eq:sum_tetb}
\end{align}
The torsion sector of the affine connection is the necessary
carrier of the TetA/TetB distinction: a torsion-free connection,
built from the parity-invariant metric alone, cannot distinguish
the two chambers. General relativity's zero-torsion condition
removes this carrier. The removed fraction is
$\Delta\mu = 1/(x+1)$, and the resulting dark-to-baryon
ratio is $1/x$. The torsion-sector field fraction and the
cold-dark-matter residual are algebraically identical under the
exchange-symmetric partition.

% ==================================================================

\begin{thebibliography}{99}

\bibitem{Nelson2026a}
S.~Nelson,
\textit{The Coordinate-Space Geometry of the Tetrad Frame},
Class.\ Quantum Grav.\ (2026), submitted (CQG-115428).

\bibitem{Hartshorne1977}
R.~Hartshorne,
\textit{Algebraic Geometry}
(Springer, New York, 1977).

\end{thebibliography}

\end{document}
